
\subsubsection{Spurious Correlation Mitigation}

Group Distributionally Robust Optimization (Sagawa et al., 2020) minimizes worst-group loss, achieving 75-80\% worst-group accuracy on Waterbirds compared to ERM's 60-75\%, but requires group annotations during training. Invariant Risk Minimization (Arjovsky et al., 2019) learns environment-invariant representations. Just Train Twice (Liu et al., 2021) identifies error-prone samples through initial training and upweights them during retraining. These methods address spurious correlations but do not explain the optimization dynamics that cause standard training to exploit shortcuts.

\subsubsection{Loss Landscape Analysis}

Li et al. (2018) introduced filter normalization for visualizing neural network loss surfaces, showing that ResNets with skip connections produce landscapes with flat, well-connected minima. Sagun et al. (2017) analyzed Hessian eigenvalue spectra using Marchenko-Pastur random matrix theory, demonstrating that outlier eigenvalues correspond to data structure. Garipov et al. (2018) showed that independently trained networks are connected by low-loss paths, enabling Fast Geometric Ensembling. These studies establish methods for analyzing loss landscape geometry but have not been applied to spurious correlation problems.

\subsubsection{Sharpness and Generalization}

Sharpness-Aware Minimization (Foret et al., 2020) seeks flat minima by perturbing weights in the direction of maximum loss increase, improving generalization. Keskar et al. (2016) demonstrated that sharp minima correlate with poor generalization. These approaches optimize for scalar flatness—minimizing maximum eigenvalues—without considering curvature orientation relative to subgroup gradients.
